%
% 11-wege.tex -- Wege in $\mathbb{C}$
%
% (c) 2026 Prof Dr Andreas Müller
%
\subsection{Wege in $\mathbb{R}^n$ und $\mathbb{C}$}
Ein \emph{Weg} in einer offenen Menge $U\subset \mathbb{R}^n$ ist eine
\index{Weg}
stetige Abbildung
\[
\gamma : I \to U : t \mapsto \gamma(t)
\]
für ein Intervall $I=[a,b]$.
Der Weg heisst \emph{geschlossen}, wenn Anfangs- und Endpunkt übereinstimmen,
\index{Weg!geschlossen}%
\index{geschlossener Weg}%
wenn also $\gamma(a)=\gamma(b)$ gilt.
Ein Weg heisst \emph{differenzierbarer Weg}, wenn die Funktion $\gamma$
\index{Weg!differenzierbar}
differenzierbar ist.
Die Ableitung
\[
\dot{\gamma}
\colon
I
\to
\mathbb{R}^n
:
t
\mapsto
\dot{\gamma}(t)
=
\frac{d\gamma}{dt}(t)
\]
heisst auch der Tangentialvektor des Weges im Punkt $\gamma(t)$.

%
% Parametertransformation auf Standardintervalle
%
\subsubsection{Parametertransformation auf Standardintervalle}
Manchmal ist es nützlich, spezielle Parameterintervalle zu verwenden.
Die Parametertransformation
\begin{equation}
t
\colon
[0,1] \to [a,b]
:
s
\mapsto
a + s (b-a)
=
(1-s)a
+
sb
\label{buch:integration:wegintegral:eqn:transformation}
\end{equation}
transformiert den Weg $\gamma\colon[a,b]\to \mathbb{R}^n$
auf den Weg $\gamma\circ t\colon [0,1]\to\mathbb{R}^n$.
Der Tangentialvektor ist nach der Kettenregel
\[
\frac{d}{ds}(\gamma\circ t)(s)
=
\dot{\gamma}(t(s))\cdot \frac{dt}{ds}
=
\dot{\gamma}(t(s))\cdot (b-a).
\]
Er wird also mit dem gleichen Faktor $b-a$ skaliert wie der Parameter.

Mit der Transformation
\eqref{buch:integration:wegintegral:eqn:transformation}
kann jeder Wege auf das Standardintervall $[0,1]$ transformiert werden.
Durch Zusammensetzung solcher Transformationen und ihrere Inversen
kann als ein Weg mit jedem beliebigen Parameterintervall
parametrisiert werden.
Wir können daher im Folgenden jederzeit annehmen, dass das Parameterintervall
das Einheitsintervall oder auch jedes beliebige andere Intervall ist.
Wir werden daher meistens das Intervall nicht im Detail spezifizieren
und es dem Leser überlassen, das Intervall aus dem Kontext abzuleiten,
sofern die Wahl des Intervalls überhaupt von Bedeutung ist.

%
% Differenzierbare Wege in C^2
%
\subsubsection{Differenzierbare Wege in $\mathbb{C}^2$}
Da $\mathbb{C}$ als zweidimensionaler reeller Vektorraum betrachtet
werden kann, ist ein Weg in $U\subset\mathbb{C}$ eine Abbildung
von einem Intervall $I$ nach $U$:
\[
\gamma
\colon
I
\to
U
:
t
\mapsto
\gamma(t).
\]
Die Ableitung nach $t$ ist ebenfalls eine komplexe Zahl, die Funktion
\[
\dot{\gamma}
\colon
I
\to
\mathbb{C}
:
t
\mapsto
\dot{\gamma}(t)
=
\frac{d\gamma}{dt}(t)
\]
heisst wieder Tangentialvektor, er ist jetzt aber auch eine komplexe 
Zahl.

%
% Real- und Imaginärteil des Weges
%
\subsubsection{Real- und Imaginärteil des Weges}
Ein Weg in $\mathbb{R}^2$ entsteht dadurch, dass man den Vektor aus
Real- und Imaginärteil von $\gamma(t)$ bildet:
\[
\gamma_{\mathbb{R}}
\colon
I
\to
\mathbb{R}^2
:
t
\mapsto
\begin{pmatrix}
\operatorname{Re}\gamma(t)\\
\operatorname{Im}\gamma(t)
\end{pmatrix}.
\]
Der Tangentialvektor besteht aus den Ableitungen von Real-
und Imaginärteil
\[
\dot{\gamma}
\colon
I
\to
\mathbb{C}
:
t
\mapsto
\frac{d}{dt}\bigl(\operatorname{Re}\gamma(t)\bigr)
+
i
\frac{d}{dt}\bigl(\operatorname{Im}\gamma(t)\bigr)
\]
nach $t$.

\input{chapters/030-integration/fig/fig-wege.tex}%

\begin{beispiel}
\label{buch:integration:wege:bsp:strecke}
Sei $U=\mathbb{C}$ und seien zwei Punkte $z_0,z_1\in \mathbb{C}$ gegeben.
Dann ist
\[
\gamma
\colon
[0,1]
\to
\mathbb{C}
:
(1-t)z_0 + tz_1
\]
ein differenzierbarer Weg, der $\gamma(0)=z_0$ mit $\gamma(1)=z_1$
verbindet.
Die Ableitung von $\gamma$ ist
\[
\dot{\gamma}(t)
=
z_1-z_0.
\]
Insbesondere ist die Ableitung konstant.

Weitere Wege mit den gleichen Endpunkten können konstruiert werden,
indem ein Weg addiert wird, der von $0$ nach $0$ führt.
Zum Beispiel ist für jedes $n\in\mathbb{N}$ 
\[
\delta(t)
=
r_0
(e^{2\pi int} - 1)
\]
ein kreisförmiger Weg mit Radius $r_0$, der $n$ Mal durchlaufen wird.
Fügt man diesen Weg und zusätzlich den Summanden $i\sin \pi t$ zum Weg
$\gamma(t)$ hinzu, erhält man den Weg
\[
\gamma_1(t)
=
\gamma(t)
+
r_0
(e^{2\pi int} - 1)
+
i\sin \pi t,
\]
der in Abbildung~\ref{buch:integration:wege:fig:wege}~a)
{\color{darkgreen}grün} eingezeichnet ist.
Der zugehörige Geschwindigkeitsvektor ist
\[
\dot{\gamma}_1(t)
=
(z_1-z_0)
+
2\pi i
n
r_0
e^{2\pi int}
+
i\pi\cos \pi t.
\qedhere
\]
\end{beispiel}

\begin{beispiel}
\label{buch:integration:wege:bsp:kreis}
Sei $U=\mathbb{C}$ und $r\in\mathbb{R}$.
Der Weg
\[
\gamma
\colon
[0,2\pi]
\to
U
:
t
\mapsto
r(\cos t + i\sin t)
\]
ist ein differenzierbarer Weg mit $\gamma(0) = \gamma(2\pi) = 1$.
Die Ableitung ist
\[
\dot{\gamma}(t)
=
r(-\sin t + i\cos t).
\]
In Satz~\ref{buch:holomorph:holomorph:satz:eulerformel}
wurde gezeigt, dass $\gamma(t) = re^{it}$ ist mit der Ableitung
$\dot{\gamma}(t)=rie^{it}$.
\end{beispiel}
